%%%%%%%%%%%%%%%%%%%%%%%%%%%%%%%%%%%%%%%%%%%%%%%%%%%%%%%%%%%%%%%%%%%%%%
% Lớp 11 | Đạo hàm | Các quy tắc đạo hàm
% Dạng 1D7?2-7: Dùng đạo hàm cho nhị thức Newton
% Số câu: 1  (sinh tự động bởi tach_cau.py; từ nay thầy sửa trực tiếp file này)
%%%%%%%%%%%%%%%%%%%%%%%%%%%%%%%%%%%%%%%%%%%%%%%%%%%%%%%%%%%%%%%%%%%%%%

%%%%%%%%%%%%%%%% TRẢ LỜI NGẮN %%%%%%%%%%%%%%%%

% nguồn: Ok_TraLoiNgan_Lop11_Lan1
\begin{ex}%[1D7V2-7]%[ma:be22148a]
	Cho biểu thức $P(x)=(x-1)^{2024}=a_0+a_1x+a_2x^2+\cdots+a_{2024}x^{2024}$. Tính giá  trị của biểu thức $$S=a_1+2a_2+3a_3+\cdots+2024a_{2024}.$$
	\shortans{0} 
	\loigiai{ 
		Đạo hàm hai  vế của biểu thức $P(x)$, ta được
		\[2024(x-1)^{2023}=a_1+2a_2x+\cdots+2024a_{2024}x^{2023}.\quad  (1)\]
		Trong $(1)$, thay $x=1$ vào hai vế  ta được
		 \[0=a_1+2a_2+3a_3+\cdots+2024a_{2024}.\]
		Vậy $S=0$. 
	} 
\end{ex}
